\documentclass[12pt]{article}
\usepackage{latexblog}

% ============================================================
% Blog Metadata
% ============================================================
\blogtitle{Java中的时间}
\blogdate{2019-11-07}
\blogtags{编程语言, 闲话编程}
\bloglang{zh}
\blogsummary{}

\begin{document}
\makeblogtitle

你如果以为，Java中谈到时间仅仅就意味着\texttt{java.util.Date}那就大错特错了，Java中的时间其实可以说五花八门，Java8发布后又增加了一些新的用来表示日期和时间的类，那么我们在构建应用程序的时候到底应该用哪个类来呢？彼此之间又有什么区别？

\section{关于时间的表示}\label{ux5173ux4e8eux65f6ux95f4ux7684ux8868ux793a}

通常由于文化和地区的不同，世界上各个地方的人们对于时间的表达方式都不尽相同，比如在中国以前用农历和十二时辰来表示，而在西方是二十四小时制（巧合的是正好一个时辰能对应上两个小时，据称与''高合成数``有关）。那么在计算机领域，在表示时间的时候也有不同的表示方法，比较常见的有：

\begin{itemize}
\tightlist
\item
  通过计算当前时间到Jan 01 1970(Unix Epoch)这一天（准确的说是00：00）经过的秒数来表示，例如 \texttt{1573090869}。
\item
  比较直接的方法就是直接以时分秒的形式表示当地时间，同时将当地的时区加进去，例如 \texttt{2001-07-04T12:08:56.235-07:00}
\end{itemize}

\section{Java中的时间类}\label{javaux4e2dux7684ux65f6ux95f4ux7c7b}

抛开时间戳不谈，在Java中专门用来表示时间的其他类有：

{\def\LTcaptype{none} % do not increment counter
\begin{longtable}[]{@{}lll@{}}
\toprule\noalign{}
Class & Since & Description \\
\midrule\noalign{}
\endhead
\bottomrule\noalign{}
\endlastfoot
java.util.Date & JDK1.0 & 日期+时间 \\
java.util.Calendar & JDK1.1 & \\
java.sql.Date & & 只包含日期 \\
java.sql.Time & & 只包含时间 \\
java.sql.Timestamp & & \\
java.time.Instant & JDK1.8 & \\
java.time.LocalTime & JDK1.8 & \\
java.time.LocalDate & JDK1.8 & \\
java.time.OffsetTime & JDK1.8 & \\
java.time.OffsetDateTime & JDK1.8 & \\
java.time.ZonedDateTime & JDK1.8 & \\
\end{longtable}
}

在Java8之前，我们通常用\texttt{java.util.Date}来表示时间，虽然没啥需求实现不了的，但有着以下的问题：

\begin{itemize}
\tightlist
\item
  时间类不统一，\texttt{java.util}和\texttt{java.sql}包中都有关于时间的类，而时间格式化的又在\texttt{java.text}包中，有点乱的很
\item
  所有时间的类都是mutable的，非线程安全
\end{itemize}

所以Java8开始对时间进行了修改，使用起来将更加方便。通常来讲，对于需要处理时区问题的系统，\texttt{OffsetDateTime}是一个较好的选择，即包含了时间信息，又包含了时区的信息，可以得到准确的时间表述。官方文档中也建议使用：

\begin{quote}
It is intended that ZonedDateTime or Instant is used to model data in simpler applications. This class may be used when modeling date-time concepts in more detail, or when communicating to a database or in a network protocol.
\end{quote}

在处理类似时间转换的时候，可以借助\texttt{ZonedDateTime}来实现。例如有一个\texttt{yyyyMMddHHmmss}格式的时间，但该时间是\texttt{CST}时间，这个时间有如下的特点：

\begin{itemize}
\tightlist
\item
  在夏季时间相当于utc+5
\item
  在冬季相当于utc+6
\end{itemize}

而每年会有一个时间点去切换这个夏季时间和冬季时间（称之为day-saving)，而且不是一个固定的日期（大致相当于按照第几个星期几来算的），要想把这个时间转换为标准的时间描述，可以采取这样的方式:

\begin{Shaded}
\begin{Highlighting}[]
\NormalTok{LocalDateTime localDateTime }\OperatorTok{=}\NormalTok{ LocalDateTime}\OperatorTok{.}\FunctionTok{parse}\OperatorTok{(}
    \StringTok{"20191010095425"}\OperatorTok{,}\NormalTok{ DateTimeFormatter}\OperatorTok{.}\FunctionTok{ofPattern}\OperatorTok{(}\StringTok{"yyyyMMddHHmmss"}\OperatorTok{));}
\NormalTok{ZonedDateTime zonedDateTime }\OperatorTok{=}\NormalTok{ ZonedDateTime}\OperatorTok{.}\FunctionTok{of}\OperatorTok{(}
\NormalTok{    localDateTime}\OperatorTok{,} \BuiltInTok{TimeZone}\OperatorTok{.}\FunctionTok{getTimeZone}\OperatorTok{(}\StringTok{"CST"}\OperatorTok{));}
\NormalTok{OffsetDateTime result }\OperatorTok{=}\NormalTok{ zonedDateTime}
    \OperatorTok{.}\FunctionTok{toOffsetDateTime}\OperatorTok{()}
    \OperatorTok{.}\FunctionTok{withOffsetSameInstant}\OperatorTok{(}\NormalTok{ZoneOffset}\OperatorTok{.}\FunctionTok{UTC}\OperatorTok{);}
\end{Highlighting}
\end{Shaded}

\begin{itemize}
\tightlist
\item
  \href{https://stackoverflow.com/questions/30234594/whats-the-difference-between-java-8-zoneddatetime-and-offsetdatetime}{What's the difference between java 8 ZonedDateTime and OffsetDateTime?}
\item
  \href{https://docs.oracle.com/javase/8/docs/api/java/time/OffsetDateTime.html}{OffsetDateTime}
\item
  \href{https://docs.oracle.com/javase/8/docs/api/java/time/ZonedDateTime.html}{ZonedDateTime}
\end{itemize}


\end{document}
